\pdfbookmark[0]{Anzeigen. Untersuchungen über Themen der höheren Geodäsie. 1846.}{anzeigen-geodaesie-zweite}
\begin{center}
\rule{.85\textwidth}{.3pt}\par
Nachrichten von der K. Gesellschaft der Wissenschaften zu Göttingen.\par 28. September 1846.\par
\rule{.85\textwidth}{.3pt}
\end{center}
\hspace*{\parindent}Am 1. September überreichte der Geheime Hofrat \textsc{Gauss} der Königlichen Gesellschaft der Wissenschaften einen Vortrag mit dem Titel
\begin{center}\textit{Untersuchungen über Themen der höheren Geodäsie, zweite Abhandlung},\end{center}
\noindent über dessen Inhalt und Zusammenhang mit der ersten Abhandlung hier kurz berichtet werden soll.

In der ersten Abhandlung wurde eine neue Methode zur Behandlung geodätischer Messungen vorgestellt. Ihr Hauptmerkmal besteht darin, dass die meisten Rechnungen ganz oder fast so geführt werden, als läge das Dreiecksnetz auf einer Kugeloberfläche statt auf einem Sphäroid, und zwar ohne jede Einbuße bei der höchsten Genauigkeit der Ergebnisse. Eine der Hauptaufgaben der geodätischen Rechnung besteht darin, aus der Länge einer Dreiecksseite, die eine geodätische Linie ist, der Breite eines Endpunkts und dem Azimut, unter dem dort der andere Endpunkt erscheint, die Breite dieses anderen Endpunkts, das dortige Azimut der Dreiecksseite und den Längenunterschied der beiden Punkte abzuleiten. Bei jener Methode beschränkt sich diese Aufgabe darauf, ein sphärisches Dreieck aufzulösen. Dennoch sind ihr in der genannten Abhandlung ein paar Seiten gewidmet, weil die gewöhnlichen Formeln der sphärischen Trigonometrie ohne Logarithmen mit mehr Stellen den Ergebnissen nicht immer eine ganz ausreichende Genauigkeit geben würden. Deshalb sind gewisse Umformungen dieser Formeln nötig. Außerdem können die Seiten von Dreiecken, deren Winkel tatsächlich gemessen werden, im Vergleich zu den Abmessungen des %
